\chapter{Start, Compile, Change}
\label{chap:start}

\begin{learningoutcomes}
  \item Compile a complete document.
  \item Recognise the preamble and document body.
  \item Type reserved characters correctly.
  \item Use the log to repair the first common errors.
\end{learningoutcomes}

\section{Build this first}

Create a folder, place a file named \file{main.tex} inside it, and type:

\begin{latexcode}{Complete example: first document}
\documentclass{article}

\begin{document}
Hello, world!

This is my first \LaTeX{} document.
\end{document}
\end{latexcode}

Compile from an editor or from a terminal:

\begin{terminalcode}{Compile and rebuild when the source changes}
latexmk -pdf main.tex
\end{terminalcode}

The result is \file{main.pdf}. Edit the words, save, and compile again. That
write--compile--inspect loop is the central skill in this book.

\section{Read the source by jobs}

\begin{center}
\begin{tabularx}{\linewidth}{@{}lX@{}}
\toprule
Source & Job \\
\midrule
\code{\textbackslash documentclass\{article\}} & Select the basic document
shape. \\
Preamble & Load packages and define reusable commands. \\
\code{\textbackslash begin\{document\}} & Start material that will be printed. \\
Body & Hold headings, paragraphs, equations, tables, and figures. \\
\code{\textbackslash end\{document\}} & Stop the document. \\
\bottomrule
\end{tabularx}
\end{center}

Commands begin with a backslash. Required arguments use braces. Optional
arguments use square brackets:

\begin{latexcode}{Command pattern}
\command[optional setting]{required content}
\end{latexcode}

\section{Add a title and sections}

\begin{latexcode}{Complete example: a short report}
\documentclass{article}

\title{A Short Measurement Report}
\author{Your Name}
\date{\today}

\begin{document}
\maketitle

\section{Question}
Do the two measurements agree?

\section{Method}
We measured each sample twice.

\section{Result}
The mean difference was small.

\section{Conclusion}
The methods produced similar values in this trial.
\end{document}
\end{latexcode}

Change the title and add a \cmd{subsection}. Do not type section numbers;
\LaTeX{} creates them.

\section{Characters that already have jobs}

The characters below are instructions in source. Add a backslash when you want
most of them printed:

\begin{latexcode}{Complete example: reserved characters}
\documentclass{article}
\begin{document}
Success rate: 95\%.

Budget: \$500.

Variables: A\_1 and A\_2.

Groups: research \& teaching.

Literal symbols: \# \{ \}.

A backslash: \textbackslash
\end{document}
\end{latexcode}

\begin{commonerror}
\textbf{Symptom:} the build stops near an ampersand or underscore.

\textbf{Fix:} use \code{\textbackslash\&} in ordinary text and
\code{\textbackslash\_} for a literal underscore. Read the first error in the
log, repair it, and compile again.
\end{commonerror}

\section{Paragraphs, spaces, and comments}

A blank source line starts a paragraph. Several ordinary spaces become one.
A percent sign starts a comment:

\begin{latexcode}{Paragraphs and comments}
This is one paragraph.
It continues on the next source line.

This is a new paragraph. % This note is not printed.
\end{latexcode}

\begin{codelab}
Compile the short report. Add one subsection, one percentage, one literal
underscore, and one comment. Then remove a closing brace, compile, read the
first error, restore the brace, and compile successfully.
\end{codelab}

\section{Keep this working pattern}

\begin{latexcode}{Minimal reusable starting file}
\documentclass{article}

% Packages and commands go here.

\title{Document Title}
\author{Author Name}
\date{\today}

\begin{document}
\maketitle

\section{Introduction}
Start writing here.

\end{document}
\end{latexcode}

\begin{chapterchecklist}
  \item The main file has a document class and matching document environment.
  \item Every opening brace or environment has a closing partner.
  \item The PDF was inspected after the latest compilation.
  \item The source and PDF are stored together.
\end{chapterchecklist}
